\uuid{rafr}
\exo7id{4680}
\titre{exo7 4680}
\auteur{quercia}
\organisation{exo7}
\datecreate{2010-03-16}
\isIndication{false}
\isCorrection{true}
\chapitre{Suite}
\sousChapitre{Convergence}
\cpge{mpsi-mp2i.nombres-reels-suites.c, pcsi.nombres-reels-suites.b, ptsi.nombres-reels-suites.b, mpsi-mp2i.nombres-reels-suites.d, pcsi.nombres-reels-suites.c, ptsi.nombres-reels-suites.c, mpsi-mp2i.nombres-reels-suites.e, pcsi.nombres-reels-suites.d, ptsi.nombres-reels-suites.d, bcpst1.suites-reelles}
\module{Analyse}
\niveau{L1}
\difficulte{}

\contenu{
\texte{
On pose $u_n = \frac {1\times 3\times 5 \times \dots \times (2n-1)}
                      {2\times 4\times 6 \times \dots \times (2n)}$.
}
\begin{enumerate}
    \item \question{Exprimer $u_n$ {\`a} l'aide de factorielles.}
    \item \question{Montrer que la suite $(u_n)$ est convergente.}
    \item \question{Soit $v_n = (n+1)u_n^2$. Montrer que la suite $(v_n)$ converge.
    Que pouvez-vous en d{\'e}duire pour $\lim_{n\to\infty}u_n$ ?}
    \item \question{On note $\alpha = \lim_{n\to\infty}v_n$. En {\'e}tudiant la suite $(nu_n^2)$,
    montrer que $\alpha > 0$.}
\reponse{
$u_n = \frac {(2n)!}{4^n(n!)^2}$.
}
\end{enumerate}
}
